\documentclass[11pt,a4paper]{article}

% ========================
% ESSENTIAL PACKAGES
% ========================

% Language and encoding
\usepackage[utf8]{inputenc}
\usepackage[T1]{fontenc}
\usepackage[english]{babel}
\usepackage{natbib}

% Math packages
\usepackage{amsmath}        % Enhanced math environments
\usepackage{amssymb}        % Additional math symbols
\usepackage{amsthm}         % Theorem environments
\usepackage{amsfonts}       % Additional math fonts
\usepackage{mathtools}      % Extension to amsmath
\usepackage{bm}             % Bold math symbols
\usepackage{dsfont}         % Double-struck fonts (for sets like ℝ, ℕ)
\usepackage{mathrsfs}       % Script fonts for math

% Graphics and figures
\usepackage{graphicx}       % Include graphics
\usepackage{float}          % Better float placement
\usepackage{subcaption}     % Subfigures
\usepackage{tikz}           % Drawing diagrams
\usepackage{pgfplots}       % Plotting
\pgfplotsset{compat=1.18}

% Page layout and formatting
\usepackage[margin=1in]{geometry}
\usepackage{fancyhdr}       % Custom headers and footers
\usepackage{titlesec}       % Customize section titles
\usepackage{enumitem}       % Customize lists
\usepackage{parskip}        % Paragraph spacing

% Colors and highlighting
\usepackage{xcolor}
\usepackage{mdframed}       % Framed environments
\usepackage{tcolorbox}      % Colored boxes

% Tables
\usepackage{array}          % Extended array environments
\usepackage{booktabs}       % Professional tables
\usepackage{multirow}       % Multi-row cells

% References and links
\usepackage{hyperref}       % Hyperlinks
\usepackage{cleveref}       % Smart cross-references

% Additional useful packages
\usepackage{cancel}         % Cancel terms in equations
\usepackage{siunitx}        % SI units
\usepackage{listings}       % Code listings (if needed)

% ========================
% CUSTOM COMMANDS
% ========================

% Common math sets
\newcommand{\R}{\mathbb{R}}
\newcommand{\E}{\mathbb{E}}
\newcommand{\N}{\mathbb{N}}
\newcommand{\Z}{\mathbb{Z}}
\newcommand{\Q}{\mathbb{Q}}
\newcommand{\C}{\mathbb{C}}

% Common operators
\DeclareMathOperator{\lcm}{lcm}
\DeclareMathOperator{\Tr}{Tr}
\DeclareMathOperator{\rank}{rank}

%\DeclareMathOperator{\span}{span}
%\DeclareMathOperator{\null}{null}

% Shortcuts for common expressions
\newcommand{\abs}[1]{\left|#1\right|}
\newcommand{\norm}[1]{\left\|#1\right\|}
\newcommand{\inner}[2]{\left\langle #1, #2 \right\rangle}
\newcommand{\ceil}[1]{\left\lceil #1 \right\rceil}
\newcommand{\floor}[1]{\left\lfloor #1 \right\rfloor}

% Bold symbols
\newcommand{\bLambda}{{\boldsymbol \Lambda}}
% Derivatives and integrals
\newcommand{\dd}[1]{\,\mathrm{d}#1}
\newcommand{\dv}[2]{\frac{\mathrm{d}#1}{\mathrm{d}#2}}
\newcommand{\pdv}[2]{\frac{\partial #1}{\partial #2}}

\renewcommand{\hat}{\widehat}
\renewcommand{\tilde}{\widetilde}

% ========================
% THEOREM ENVIRONMENTS
% ========================

\theoremstyle{definition}
\newtheorem{definition}{Definition}[section]
\newtheorem{example}{Example}[section]
\newtheorem{exercise}{Exercise}[section]

\theoremstyle{plain}
\newtheorem{theorem}{Theorem}[section]
\newtheorem{lemma}{Lemma}[section]
\newtheorem{corollary}{Corollary}[section]
\newtheorem{proposition}{Proposition}[section]

\theoremstyle{remark}
\newtheorem{remark}{Remark}[section]
\newtheorem{note}{Note}[section]

\newcommand{\supp}{\mathrm{supp}}
\newcommand{\op}{\mathrm{op}}

% ========================
% COLORED BOXES
% ========================

% Important theorem box
\newtcolorbox{importantbox}{
	colback=blue!5!white,
	colframe=blue!75!black,
	title=Important
}

% Warning/Note box
\newtcolorbox{notebox}{
	colback=yellow!10!white,
	colframe=orange!75!black,
	title=Note
}

% Example box
\newtcolorbox{examplebox}{
	colback=green!5!white,
	colframe=green!75!black,
	title=Example
}

% ========================
% DOCUMENT SETUP
% ========================

% Page style
\pagestyle{fancy}
\fancyhf{}
\rhead{\thepage}
\lhead{\leftmark}

% Title formatting
\title{Functional Bandit Notes}
\author{Your Name}
\date{\today}

% Hyperlink setup
\hypersetup{
	colorlinks=true,
	linkcolor=blue,
	filecolor=magenta,      
	urlcolor=cyan,
	citecolor=red
}

\begin{document}

Let $x_1,\cdots,x_n\overset{i.i.d.}{\sim} P$ be square integrable stochastic process where $P$ is supported on a compact set $\mathcal T\subset \mathbb R.$ Let $y_i=\langle x_i,\beta_0 \rangle+\varepsilon_i$ (w.l.o.g. suppose $x$ is centered and we omit the intercept here). Let $\boldsymbol{K}$ be a known reproducing kernel and let $\mathcal H(\boldsymbol{K})$ be the corresponding RKHS. Assume that $\beta_0\in\mathcal H(\boldsymbol K).$ Let $\|\cdot\|_{\psi_2}$ denote the sub-gaussian norm of a random variable. Suppose $L_{\boldsymbol{K}^{1/2}}x$ is sub-gaussian (has bounded sub-gaussian norm).

Let $\boldsymbol{C}_0$ be the covariance function of $x$ defined as follows:
\[\boldsymbol{C}_0(s,t)=\mathbb{E}[(x(s)-\mathbb{E}[x(s)])(x(t)-\mathbb{E}[x(t)])].\]
 Suppose that the eigenvalues $s_k$ of the operaotr $\boldsymbol T_0=L_{\boldsymbol{K}^{1/2}\boldsymbol{C}_0\boldsymbol{K}^{1/2}}$ have $s_k\asymp k^{-2r}.$ We have the subgaussian norm $\sigma_x\asymp s_1.$ (We shall assume that $r>1/2$ such that the trace of operator is bounded).
Let also $\boldsymbol{T}_\lambda=\boldsymbol{T}_0+\lambda \bold{1}.$ Denote $\hat{\boldsymbol{T}}_n$ to be
\[\hat{\boldsymbol{T}}_n=L_{\boldsymbol{K}^{1/2}\boldsymbol{V}_n\boldsymbol{K}^{1/2}}\]
where \(\boldsymbol{V}_n(s,t)=\frac{1}{n}\sum_{i=1}^n (x_i(s)-\bar{x}(s))(x_i(t)-\bar{x}(t))\)
and let $\hat{\boldsymbol{T}}_\lambda=\hat{\boldsymbol{T}}_n+\lambda\bold{1}.$

Throughout we omit poly-logarithmic terms in $\tilde O(\cdot)$ and $\lesssim$ notations. When we say ``with high probability'', it means with probability $1-O(T^{-4})$ or higher degrees.



The objective is to construct UCB $\Delta(\cdot)$ so that it satisfies the following properties:
\begin{enumerate}
\item With high probability $|x^\top(\hat\beta-\beta_0)|\leq\Delta(x)$ for all $x\in\mathcal X$, where $\mathcal X$ contains $T$ vectors in the support of $P$;
\item $\mathbb E_P[\Delta(x)]\leq \Psi$ where $\Psi$ is on the order of $\tilde O(1/\sqrt{n})$.
\end{enumerate}

Note that we don't require $\Delta(x)$ to be valid \emph{uniformly} on the support of $P$, but it has to be valid pointwise with high probability so that a union bound of $T$ vectors won't hurt.

%\subsection{Successive elimination}
%
%It is possible to use such an UCB to do successive elimination, without requiring $\lambda_{\min}$ conditions. Algorithm operates as below:
%\begin{enumerate}
%\item For epochs $\tau=1,2,3,\cdots$ each of length $n_\tau=2^\tau$, do the following:
%\begin{enumerate}
%\item Obtain $\hat\beta_\tau^1$, $\hat\beta_\tau^2$ and $\Delta_\tau^1(\cdot)$, $\Delta_\tau^2(\cdot)$ from all samples in the previous epoch;
%\item For all $t$ in epoch $\tau$, observe $x_t\sim P$ and do the following:
%\begin{enumerate}
%\item If $x_t^\top\hat\beta_s^1-\Delta_s^1(x_t)>x_t^\top\hat\beta_s^2+\Delta_s^2(x_t)$ for any $s\leq \tau$ then take action 1, and similarly in the other direction take action 2;
%\item Otherwise, take action 1 or 2 with half probability each.
%\end{enumerate}
%\end{enumerate}
%\end{enumerate}
%
%\paragraph{Analysis.} With high probability the optimal action is never eliminated. So regret is only incurred when actions are taken with half probability each. 
%Define distribution $z\sim Q_\tau$ as the following: first sample $x\sim P$, then set $z=x$ if on $x$ actions are taken with half probability each during epoch $\tau$, and $z=0$ otherwise, with the understanding that $\Delta_\tau^1(0)=\Delta_\tau^2(0)=0$.
%Let $P_\tau^1,P_\tau^2$ be the distributions received by actions 1 and 2 from all samples in the previous epoch.
%We have the following observations:
%\begin{enumerate}
%\item $\mathbb E[\Delta_\tau^1(x)|P_\tau^1],\mathbb E[\Delta_\tau^2(x)|P_\tau^2]\lesssim \Psi_\tau \lesssim 1/\sqrt{n_\tau}$, where the last inequality is hiding $\mathrm{poly}(d)$ factors. This holds true as long as the two actions are balanced.
%\item For any $x\in\supp(Q_\tau)\backslash\{0\}$, $p(x|Q_\tau)\leq \max\{2\kappa_{\min}^{-1}p(x|P_\tau^1),2\kappa_{\min}^{-1}p(x|P_\tau^2)\}$ where $\kappa_{\min}=\min\{P(a^*(x)=1),P(a^*(x)=2)\}$ because $x$ is not eliminated during epoch $\tau-1$ and therefore has half probability being assigned to each action.
%\end{enumerate}
%
%Therefore, during epoch $\tau$, the cumulative regret is
%\begin{align*}
%n_\tau&\mathbb E[\Delta_\tau^1(z)+\Delta_\tau^2(z)|Q_\tau] \lesssim n_\tau(\mathbb E[\Delta_\tau^1(x)|P_\tau^1] + \mathbb E[\Delta_\tau^2(x)|P_\tau^2]) \lesssim n_\tau\times \frac{1}{\sqrt{n_\tau}} = \tilde O(\sqrt{n_\tau}).
%\end{align*}
%This leads to $\tilde O(\sqrt{T})$ regret. Dimension dependency depends on $\Psi$ and varies.

\subsection{Direct analysis of ridge regression when $\bold{T}$ is known}

Define the ridge regressor of $\tilde{\beta}_0 = (L_{\boldsymbol{K}^{1/2}})^{-1}\beta_0$ to be
\[\hat{\beta}_\Lambda = \hat{\boldsymbol{T}}_\Lambda^{-1}(\hat{\boldsymbol T}_n\tilde{\beta}_0+g_\epsilon)\]
where 
\[g_\epsilon = \frac{1}{n}\sum_{i=1}^n \varepsilon_iL_{\boldsymbol K^{1/2}}x_i.\]
Suppose that $\tilde{x}$ is a gaussian process and $\tilde{x}\sim \mathcal{N}(0,\Sigma),$ where $\Sigma$ is the covariance operator. Than we have with high probability
\[\langle \tilde x, \tilde\beta_0-\hat{\beta}_\bLambda \rangle\leq \]
For any $\tilde x$ we have
\[\langle \tilde x, \tilde\beta_0-\hat{\beta}_\bLambda \rangle=\langle \tilde x, \tilde{\beta}_0-\tilde{\beta}_\bLambda \rangle+\langle \tilde x, \tilde{\beta}_\bLambda-\hat{\beta}_\bLambda\rangle,\]
where
\[\tilde{\beta}_\bLambda=\boldsymbol{T}_{\bLambda}^{-1}\boldsymbol{T}\tilde{\beta_0}.\]
% \[\langle \tilde x, \tilde\beta_0-\hat{\beta}_\bLambda \rangle=\langle \tilde{x},  -\bLambda \hat{\boldsymbol T}_\bLambda^{-1}\tilde\beta_0+\hat{\boldsymbol T}_\bLambda^{-1}\hat{\boldsymbol T}_n g_\epsilon\rangle.\]

\paragraph{Analysis of the population bias term.}
We have
\begin{align*}
	\langle \tilde x, \tilde\beta_0-\hat{\beta}_\bLambda \rangle = \langle \tilde x,-\bLambda  {\boldsymbol T}_{\bLambda}^{-1}\tilde\beta_0\rangle=\sum_{j=1}^{\infty} \tilde{\theta}_j\frac{\alpha_j\lambda_j}{\lambda_j+s_j}.
\end{align*}



\paragraph{Analysis of the bias term.}
We have
\begin{align}
	\langle L_{\boldsymbol K^{1/2}}x, \lambda \hat{\boldsymbol{T}}_\lambda^{-1}f_0\rangle=\lambda 	\langle L_{\boldsymbol K^{1/2}}x,  \hat{\boldsymbol{T}}_\lambda^{-1}f_0\rangle = \lambda 	\langle L_{\boldsymbol K^{1/2}}x,  \hat{\boldsymbol{T}}_\lambda^{-1}\hat{f}_\lambda\rangle + \lambda 	\langle L_{\boldsymbol K^{1/2}}x,  \hat{\boldsymbol{T}}_\lambda^{-1}(f_0-\hat{f}_\lambda)\rangle .\label{ineq: bias}
\end{align}
For the RHS, the former term is explicit, while we do not know the second term. Thus we can write
\begin{align*}
	\lambda 	\langle L_{\boldsymbol K^{1/2}}x,  \hat{\boldsymbol{T}}_\lambda^{-1}(f_0-\hat{f}_\lambda)\rangle&\leq \lambda \left|{\langle L_{\boldsymbol K^{1/2}} x,\hat{\boldsymbol T}_{\lambda}^{-3/2} L_{\boldsymbol K^{1/2}} x\rangle}\right|\left\|\hat{\boldsymbol{T}}_{\lambda}^{1/2}(f_0-\hat{f}_{\lambda})\right\|_{\mathcal L_2}\\&\preceq\lambda \left|{\langle L_{\boldsymbol K^{1/2}} x,\hat{\boldsymbol T}_{\lambda}^{-3/2} L_{\boldsymbol K^{1/2}} x\rangle}\right|\left\|{\boldsymbol{T}}_{\lambda}^{1/2}(f_0-\hat{f}_{\lambda})\right\|_{\mathcal L_2},
\end{align*}
given $\hat{\boldsymbol T}_{\lambda}\succeq (1/2)\boldsymbol{T}_{\lambda}.$ (To obtain this, we need to use Berstein that consider the effective dimension like \cite{minsker2011some}). We further have
\[\left\|{\boldsymbol{T}}_{\lambda}^{1/2}(f_0-\hat{f}_{\lambda})\right\|_{\mathcal L_2}\leq \sqrt{\lambda}+\left\|{\boldsymbol{T}}_{0}^{1/2}(f_0-\hat{f}_{\lambda})\right\|_{\mathcal L_2}\]
and
\[\left\|{\boldsymbol{T}}_{0}^{1/2}(f_0-\hat{f}_{\lambda})\right\|_{\mathcal L_2}\preceq \sqrt{\lambda}+\left\|\boldsymbol T_0^{1/2}\boldsymbol{T}_{\lambda}^{-1/2}g_{\varepsilon}\right\|_{\mathcal L_2}\preceq \sqrt{\lambda}+\sqrt{\frac{1}{n}}\sqrt{{\langle L_{\boldsymbol{K}^{1/2}}x,\hat{\boldsymbol T}_\lambda^{-1} L_{\boldsymbol K^{1/2}}x\rangle}},\]
where the first term on the RHS is the deterministic bias term and the second term is the variance term. Putting everything together we have that the point-wise bias in \eqref{ineq: bias} can be upper bounded by
\begin{align*}
		&\langle L_{\boldsymbol K^{1/2}}x, \lambda \hat{\boldsymbol{T}}_\lambda^{-1}f_0\rangle\\
		&\preceq\lambda 	\langle L_{\boldsymbol K^{1/2}}x,  \hat{\boldsymbol{T}}_\lambda^{-1}\hat{f}_\lambda\rangle + \lambda \left|{\langle L_{\boldsymbol K^{1/2}} x,\hat{\boldsymbol T}_{\lambda}^{-3/2} L_{\boldsymbol K^{1/2}} x\rangle}\right|\left(\sqrt{\lambda}+\sqrt{\frac{1}{n}}\sqrt{{\langle L_{\boldsymbol{K}^{1/2}}x,\hat{\boldsymbol T}_\lambda^{-1} L_{\boldsymbol K^{1/2}}x\rangle}}\right)\\
		&:=\Delta_b^1(x)+\Delta_b^2(x)
		:=\Delta_b(x).
\end{align*}
\paragraph{Analysis of the variance term.}
We have the variance term
\[v=\langle L_{\boldsymbol K^{1/2}}x,\hat{\boldsymbol T}_{\lambda}^{-1}g_\epsilon\rangle=\frac{1}{n}\sum_{i=1}^n \varepsilon_i\langle L_{\boldsymbol K^{1/2}}x,\hat{\boldsymbol T}_{\lambda}^{-1}L_{\boldsymbol K^{1/2}}x_i\rangle.\]
Since $\varepsilon_i$ are i.i.d. $\sigma^2$ subgaussian, we have
\[v\preceq\sqrt{\frac{1}{n}}\sqrt{{\langle L_{\boldsymbol{K}^{1/2}}x,\hat{\boldsymbol T}_\lambda^{-2}\hat{\boldsymbol T}_n L_{\boldsymbol K^{1/2}}x\rangle}}\leq \sqrt{\frac{1}{n}}\sqrt{{\langle L_{\boldsymbol{K}^{1/2}}x,\hat{\boldsymbol T}_\lambda^{-1} L_{\boldsymbol K^{1/2}}x\rangle}}:=\Delta_v(x).\]
\paragraph{Putting bias and variance together.} 
We can now define \[\Delta(x):=\Delta_b(x)+\Delta_v(x)\]
and bound  $\mathbb E_P[\Delta(x)].$ For the bias term, we have
\begin{align*}
\mathbb{E}\left[\Delta_b^1(x)\right]\preceq\lambda\sqrt{\mathbb{E}\left[	\langle L_{\boldsymbol K^{1/2}}x,  \hat{\boldsymbol{T}}_\lambda^{-1}\hat{f}_\lambda\rangle ^2\right]}\preceq\lambda\sqrt{\mathbb{E}\left[	\langle L_{\boldsymbol K^{1/2}}x,  {\boldsymbol{T}}_\lambda^{-1}\hat{f}_\lambda\rangle ^2\right]}\preceq \lambda^{1/2}
\end{align*}
and
\begin{align*}
	\mathbb E_P\left[\Delta_b^2(x)\right]&\preceq \lambda \sqrt{\mathbb{E}_P\left[\langle L_{\boldsymbol K^{1/2}}x,\hat{\boldsymbol{T}}_{\lambda}^{-3/2} L_{\boldsymbol K^{1/2}} x\rangle^2\right]}\left(\sqrt{\lambda}+\sqrt{\frac{1}{n}}\sqrt{\mathbb{E}_P\left[{\langle L_{\boldsymbol{K}^{1/2}}x,\hat{\boldsymbol T}_\lambda^{-1} L_{\boldsymbol K^{1/2}}x\rangle}\right]}\right).
\end{align*}
For the first expectation on the RHS, we have
\begin{align}
\lambda	\sqrt{\mathbb{E}_P\left[\langle L_{\boldsymbol K^{1/2}}x,\hat{\boldsymbol{T}}_{\lambda}^{-3/2} L_{\boldsymbol K^{1/2}} x\rangle^2\right]}& \preceq\lambda \sqrt{\mathbb{E}_P\left[\langle L_{\boldsymbol K^{1/2}}x,{\boldsymbol{T}}_{\lambda}^{-3/2}\rangle^2\right]}=\sqrt{\sum_{k=1}^{\infty}\frac{s_k\lambda^2}{(\lambda+s_k)^3}}.\nonumber
	%\label{ineq: first_bias_expectation}
\end{align}
which yields
\begin{align*}
	\mathbb E_P\left[\Delta_b^2(x)\right]&\preceq \sqrt{\mathbb{E}_P\left[\langle L_{\boldsymbol K^{1/2}}x,{\boldsymbol{T}}_{\lambda}^{-1} L_{\boldsymbol K^{1/2}} x\rangle^2\right]}\left(\sqrt{\lambda}+\sqrt{\frac{1}{n}}\sqrt{\mathbb{E}_P\left[{\langle L_{\boldsymbol{K}^{1/2}}x,\hat{\boldsymbol T}_\lambda^{-1} L_{\boldsymbol K^{1/2}}x\rangle}\right]}\right)\\
	&\preceq \sqrt{\lambda}+\sqrt{\frac{1}{n}}\sqrt{\mathbb{E}_P\left[{\langle L_{\boldsymbol{K}^{1/2}}x,\hat{\boldsymbol T}_\lambda^{-1} L_{\boldsymbol K^{1/2}}x\rangle}\right]},
\end{align*}
since \(\sqrt{\mathbb{E}_P\left[\langle L_{\boldsymbol K^{1/2}}x,{\boldsymbol{T}}_{\lambda}^{-1} L_{\boldsymbol K^{1/2}} x\rangle^2\right]}=o(1).\) 
We further have
\begin{align*}
\sqrt{\mathbb E\left[{{\langle L_{\boldsymbol{K}^{1/2}}x,\hat{\boldsymbol T}_\lambda^{-1}L_{\boldsymbol K^{1/2}}x\rangle}}\right]}
	&\preceq\sqrt{\mathbb E\left[{{\langle L_{\boldsymbol{K}^{1/2}}x,{\boldsymbol T}_\lambda^{-1} L_{\boldsymbol K^{1/2}}x\rangle}}\right]}\\
	&=\sqrt{\sum_{k=1}^{\infty}\frac{s_k}{\lambda+s_k}},
\end{align*}
where $s_k$ are the eigenvalues of $\boldsymbol{T}_0$ given that $\hat{\boldsymbol T}_{\lambda}\succeq (1/2) \boldsymbol T_\lambda.$ Since $s_k\asymp k^{-2r},$ we have
\[\sqrt{\sum_{k=1}^{\infty}\frac{s_k}{\lambda+s_k}}\asymp \exp(\log \lambda/(-2r)).\]
For the variance term, we similarly have 
\[\mathbb{E}_P\left[\Delta_v(x)\right]\leq \sqrt{\frac{1}{n}\exp(\log\lambda/(-2r))}.\]
Putting everything together, we have
\[\mathbb{E}_P\left[\Delta_b(x)+\Delta_v(x)\right]\preceq\mathbb{E}_P\left[\Delta_b^1(x)\right] +\mathbb{E}_P\left[\Delta_b^2(x)\right]+\mathbb{E}_P\left[\Delta_v(x)\right]\preceq\sqrt{\lambda}+\sqrt{\frac{1}{n}\exp(\log \lambda/(-2r))}.\]
By choosing $\lambda\asymp n^{-2r/(2r+1)}$ we have 
\[b+v\asymp n^{-r/(2r+1)},\]
which is minimax optimal.


\bibliographystyle{plain}
\bibliography{refs}

\end{document}